% The next line makes it so it compiles the main file if I hit the compile button
% when editing this file in TeXworks.
% !TEX root = ../ActiveIntroToDiscreteMathAndAlgorithms.tex
%%%%%%%%%%%%%%%%%%%%%%%%%%%%%%%%%%%%
\section{Problems}

\begin{pro}
Find at least three {\em different} sequences that begin with 1, 3, 7
whose terms are generated by a simple formula or rule.  
By different, I mean none of the sequences can have exactly the same terms.  In other words,
your answer cannot simply be three different ways to generate the same sequence.
\end{pro}

\begin{pro}
Let $q_n = 2 q_{n-1} + 2n + 5$, and $q_0=0$. 
Compute $q_1$, $q_2$, $q_3$ and $q_4$.
\end{pro}

\begin{pro}
Let $a_n = a_{n-2}+n$, $a_0=0$, and $a_1=1$.  
Compute $a_2$, $a_3$, $a_4$ and $a_5$.
\end{pro}

\begin{pro}
Let $a_n = n\times a_{n-1}+5$, and $a_0=1$.  
Compute $a_1$, $a_2$, $a_3$, $a_4$ and $a_5$.
\end{pro}

\begin{pro}
Define a sequence $\{x_n\}$ by $x_0=1$, and $x_n=2x_{n-1}+1$ if $n\geq 1$.
Find a closed form for the $n$th term of this sequence.
{\em Prove that your solution is correct.}
\end{pro}


\begin{pro}
Compute each of the following:
\begin{multicols}{3}
\begin{lenum}
\item ${\displaystyle\sum_{k=5}^{40}k}$
\item ${\displaystyle\sum_{j=5}^{22}(2^{j+1} - 2^j)}$
\item ${\displaystyle\sum_{k=0}^{n}5k}$
\item ${\displaystyle\sum_{i=1}^3\sum_{j=1}^4 j}$
\item ${\displaystyle\sum_{k=1}^{n}k(k-1)}$
\item ${\displaystyle\sum_{j=1}^{ n }5^j}$
\item ${\displaystyle\sum_{j=0}^{ \log_2 n}2^j}$
\item ${\displaystyle\sum_{i=0}^{\log_2 n}
\left(\frac{n}{2^i}\right)}$
\item ${\displaystyle\sum_{i=1}^n\sum_{j=1}^i \sum_{k=1}^j 1}$
\end{lenum}
\end{multicols}
\end{pro}

\begin{pro}
Here is a standard interview question for prospective computer programmers: 
You are given a list of $1,000,001$ positive integers from the set
$\{1,2,\ldots , 1,000,000\}$. 
In the list, every member of  $\{1,2,\ldots , 1,000,000\}$ 
is listed once, except for $x$, which is listed twice, and the numbers are listed in some unknown order. 
How do you find what $x$ is without doing a $1,000,000$ step search (e.g. check if 1 is on the list twice, then check if 2 is on 
the list twice, etc.)?
How much faster is your solution than the naive solution?
\end{pro}

\begin{pro} Find a closed formula for
$$T_n = 1^2 - 2^2 + 3^2 - 4^2 + \cdots + (-1)^{n - 1}n^2.$$
%\begin{answerSer}
%$\dfrac{n(n+1)(-1)^{n+1}}{2}$.
%\end{answerSer}
\end{pro}

\begin{pro} Show that when $n\geq 1$, 
$$1 + 3 + 5 + \cdots + (2n - 1) = n^2.$$
\end{pro}

\begin{pro} Assuming $n\geq 1$, find and prove a closed formula for   
$$2+4+6+\cdots+ 2n$$
\end{pro}

\begin{pro} Show that when $n\geq 1$, 
$$\sum _{k = 1} ^n \frac{k}{k^4 + k^2 + 1} = \frac{1}{2}\cdot\frac{n^2 + n}{n^2 + n + 1}.$$
\end{pro}

\newpage

\begin{pro}
Legend says that the inventor of the game of chess, Sissa ben Dahir,  asked the King Shirham of India
to place a grain of wheat on the first square of the
chessboard, $2$ on the second square, $4$ on the third square, $8$
on the fourth square, etc..
\begin{lenum}
\item  How many grains of wheat are to be
put on the last ($64$-th) square?
\item  How many grains, total, are
needed in order to satisfy the greedy inventor?
\item  Given that $15$
grains of wheat weigh approximately one gram, what is the
approximate weight, in kg, of the wheat needed?
\item  Given that the
annual production of wheat is $350$ million tonnes, how many years,
approximately, are needed in order to satisfy the inventor (assume
that production of wheat stays constant)?
\end{lenum}
%\begin{answerSer}\noindent\begin{enumerate}
%\item
%$2^{63} = 9223372036854775808$, \item  $2^{64} - 1 =
%18446744073709551614$, \item  $1.2\times 10^{15}$ kg, or $1200$ billion
%tonnes \item  3500 years
%\end{enumerate}
%\end{answerSer}
\end{pro}

% Eval example now.
%\begin{pro}
%Consider the following function:
%\begin{code}
%int ferzle(int n) {
%    if(n<=0) {
%        return 3;
%    } else {
%        return ferzle(n-1) + 2;
%    }
%}
%\end{code}
%\begin{enumerate}
%\item
%Determine what \verb+ferzle(n)+ returns for $n=0, 1, 2, 3, 4$.
%\item
%Re-write \verb+ferzle+ without using recursion.
%\end{enumerate}
%\end{pro}



\begin{pro}
It is easy to see that we can define $n!$ recursively by defining $0!=1$, and if $n>0$,
$n!=n\cdot (n-1)!$. 
Does the following method correctly compute $n!$?  
If not, state what is wrong with it and fix it.
\begin{code}
	int factorial(int n) {
	    return n * factorial(n-1);
	}
\end{code}
\end{pro}

\begin{pro}
Find a closed formula for $\dsum _{k=1}^n k^2(k-1)$. 
Simplify the formula as much as possible.
\end{pro}

\begin{pro}
Find a closed formula for $\dsum _{k=1}^n k\cdot k!$. 
(Hint: What is $(k+1)! - k!$, and why does it matter?)
Simplify the formula as much as possible.

%\begin{answer}
%From the hint: $k\cdot k! = (k+1)! -k!$ and we get the telescoping sum
%$$\sum _{1\leq k \leq n} k\cdot k! =  
%\sum _{1\leq k \leq n} (k+1)! -k!= (2!-1!)+(3!-2!)+(4!-3!)+\cdots ((n+1)!-n!) = (n+1)!-1!.  $$
%\end{answer}
\end{pro}

\begin{pro}
Prove that for $n\geq 1$, 
$\displaystyle \dsum_{k=1}^{n} k^3 = \left(\sum_{k=1}^{n} k \right)^2$.
\end{pro}


\begin{pro}
A student turned in the code below (which does as its name suggests).  
I gave them a `C' on the assignment because although it works,
it is very inefficient. 
\begin{code}
	int sumFromOneToN(int n) {
	    int sum = 0;
	    for(int i=1;i<=n;i++) {
	        sum = sum + i;
	    }
	return sum;
	}
\end{code}
\begin{lenum}
\item
Write the `A' version of the algorithm (in other words, a more efficient version).
You can assume that $n\geq 1$.  
\item 
Compute \verb+sumFromOneToN(30)+ based on your algorithm.
\end{lenum}
\end{pro}

\newpage
%Now an eval example.  But in a later chapter.  Keep it here.
\begin{pro}
A student turned in the code below (which does as its name suggests).  
I gave them a `C' on the assignment because although it works,
it is very inefficient.
\begin{code}
	int sumFromMToN(int m, int n) {
	    int sum = 0;
	    for(int i=1;i<=n;i++) {
	        sum = sum + i;
	    }
	    for(int i=1;i<m;i++) {
	        sum = sum - i;
	    }
	    return sum;
	}
\end{code}
\begin{lenum}
\item
Write the `A' version of the algorithm (in other words, a more efficient version).
You can assume that $1\leq m\leq n$.
\item 
Compute \verb+sumFromMToN(10,50)+ based on your algorithm.
\end{lenum}
\end{pro}

\begin{pro}
How many times does the function \verb+foo+ get called in the following code?
\begin{code}
     for(int i=0;i<n;i++) {
         for(int j=0;j<i;j++) {
             foo(i,j);
         }
      }
\end{code}
\end{pro}

\begin{pro}
Consider the following code.
\begin{code}
    int g(int n) {
        int sum = 0;
        for(int i=1;i<=n;i++) {
            for(int j=1;j<=n;j++) {
                sum = sum+1;
            }
         }
         return sum;
    }
\end{code}
\begin{lenum}
\item What function does the code compute? In other words, what is $g(n)$?
\item Write a more efficient version of the code.
\end{lenum}
\end{pro}

\begin{pro}
Consider the following code, where \verb+power(2,i)+ computes $2^i$.
\begin{code}
    int h(int n) {
        int sum = 0;
        for(int i=0;i<n;i++) {
            sum = sum+power(2,i);
         }
         return sum + 1;
    }
\end{code}
\begin{lenum}
\item What function does the code compute? In other words, what is $h(n)$?
\item Write a more efficient version of the code.
\end{lenum}
\end{pro}



\begin{pro}
Write out explicitly the $3\times 3$ matrix $A = [a_{ij}]$ where
$a_{ij} = ij$.
%\begin{answer} $A = \begin{bmatrix} 1 & 2 & 3 \cr
%2 & 4 & 6 \cr 3 & 6 & 9 \cr \end{bmatrix}.$  
%\end{answer}
\end{pro}


\begin{pro}
Determine $2\times 2$ matrices $A$ and $B$ such that
$$2A - 5B = \begin{bmatrix} 1 & -2 \cr 0 & 1 \cr \end{bmatrix} \mbox{ and } -2A + 6B = \begin{bmatrix} 4 & 2 \cr 6 & 0 \cr\end{bmatrix}.   $$
%\begin{answer} $A = \begin{bmatrix}
%13 & -1 \cr 15 & 3 \cr \end{bmatrix}$, 
%$B = \begin{bmatrix}5 & 0 \cr 6 & 1 \cr \end{bmatrix}$ %\end{answer}
\end{pro}

%\begin{pro}
%Let $A=[a_{ij}]\in\mat{n\times n}{\BBR}$. Prove that
%$$\min_j\max_ia_{ij}\geq\max_i\min_ja_{ij}.$$
%\end{pro}

\begin{pro}
A person goes along the rows of a movie theater and asks the
tallest person of each row to stand up. Then he selects the
shortest of these people, who we will call the {\em shortest
giant}. Another person goes along the rows and asks the shortest
person to stand up and from these he selects the tallest, which we
will call the {\em tallest dwarf}. Who is taller, the tallest
dwarf or the shortest giant? Prove your answer.
\end{pro}

\begin{pro}[Putnam Exam, 1959]
Choose five elements from the following matrix, 
no two coming from the same row or column, so that
the minimum of these five elements is as large as possible.
$$\begin{bmatrix}11 & 17 & 25
& 19 & 16 \cr 24 & 10 & 13 & 15 & 3 \cr 12 & 5 & 14 & 2 & 18 \cr
23 & 4 & 1 & 8 & 22 \cr 6 & 20 & 7 & 21 & 9
\end{bmatrix}$$ 
%\begin{answer}
%The set of border elements is the union of two rows and two columns.
%Thus we may choose at most four elements from the border, and at
%least one from the central $3\times 3$ matrix. The largest element
%of this $3\times 3$ matrix is $15$, so any allowable choice of does
%not exceed $15$. The choice $25$, $15$, $18$,l $23$, $20$ shews that
%the largest minimum is indeed $15$.
%\end{answer}
\end{pro}

\begin{pro}
Find $a+b+c$ if $\begin{bmatrix}1 & 2 & 3 \cr 2 & 3 & 1 \cr 3 & 1 &
2 \cr \end{bmatrix}
\begin{bmatrix}1 & 1 & 1 \cr 2 & 2 & 2 \cr 3 & 3 & 3 \cr \end{bmatrix} =
\begin{bmatrix}a & a & a \cr b & b & b \cr c & c & c \cr \end{bmatrix}$.
%\begin{answer}
%$\begin{bmatrix}1 & 2 & 3 \cr 2 & 3 & 1 \cr 3 & 1 & 2 \cr
%\end{bmatrix}
%\begin{bmatrix}1 & 1 & 1 \cr 2 & 2 & 2 \cr 3 & 3 & 3 \cr \end{bmatrix} =
%\begin{bmatrix}14 & 14 & 14 \cr 11 & 11 & 11 \cr 11 & 11 & 11 \cr \end{bmatrix}$, whence $a+b+c=36$.
%\end{answer}
\end{pro}


%\begin{pro}
%Let
%$$ A = \begin{bmatrix}
%\overline{2} & \overline{3} & \overline{4} & \overline{1} \cr
% \overline{1} & \overline{2} & \overline{3} & \overline{4} \cr
%  \overline{4} & \overline{1} & \overline{2} & \overline{3} \cr
%   \overline{3} & \overline{4} & \overline{1} & \overline{2} \cr
%\end{bmatrix}, \ \ \ B = \begin{bmatrix}
%\overline{1} & \overline{1} & \overline{1} & \overline{1} \cr
%\overline{1} & \overline{1} & \overline{1} & \overline{1} \cr
%\overline{1} & \overline{1} & \overline{1} & \overline{1} \cr
%\overline{1} & \overline{1} & \overline{1} & \overline{1} \cr
%\end{bmatrix}$$ be matrices in $\mat{4\times 4}{\BBZ_5}$ . Find the
%products $AB$ and $BA$.
%\begin{answer}
%$AB = {\bf 0}_4$ and $ BA = \begin{bmatrix} \overline{0} &
%\overline{2} & \overline{0} & \overline{3} \cr \overline{0} &
%\overline{2} & \overline{0} & \overline{3} \cr \overline{0} &
%\overline{2} & \overline{0} & \overline{3} \cr \overline{0} &
%\overline{2} & \overline{0} & \overline{3} \cr
%\end{bmatrix}.$
%\end{answer}
%\end{pro}


\begin{pro}
(Requires calculus)
Let $$ A = \begin{bmatrix} 0 & \dfrac{1}{2} & 0 \cr \dfrac{1}{2}  &
0 & 0 \cr 0 & 0 & \dfrac{1}{2} \cr \end{bmatrix}. $$ Calculate the
value of the infinite series
$${\bf I}_3 + A + A^2 +A^3 + \cdots .  $$
%\begin{answer}
%For this problem you need to recall that if $|r|<1$, then
%$$a+ar + ar^2+ar^3+\cdots + \cdots =\dfrac{a}{1-r}.  $$
%This gives
%$$ 1 + \tfrac{1}{4} +  \tfrac{1}{4^2}+  \tfrac{1}{4^3}+\cdots    = \dfrac{1}{1-\tfrac{1}{4}} = \dfrac{4}{3}, $$
%$$\tfrac{1}{2} +  \tfrac{1}{2^3}+ \tfrac{1}{2^5}+\cdots  = \dfrac{\tfrac{1}{2}}{1-\tfrac{1}{4}} = \dfrac{2}{3},  $$
%and
%$$1 + \tfrac{1}{2} +  \tfrac{1}{2^2}+  \tfrac{1}{2^3}+\cdots  = \dfrac{1}{1-\tfrac{1}{2}}  = 2.  $$
% By looking at a few small cases, it is easy to establish by
%induction that for $n\geq 1$
%$$ A^{2n-1}= \begin{bmatrix} 0 & \dfrac{1}{2^{2n-1}} & 0 \cr \dfrac{1}{2^{2n-1}}  & 0 & 0 \cr 0 & 0 & \dfrac{1}{2^{2n-1}} \cr \end{bmatrix},
%\qquad  A^{2n} = \begin{bmatrix}\dfrac{1}{2^{2n}} & 0  & 0 \cr 0 &
%\dfrac{1}{2^{2n}} & 0 \cr 0 & 0 & \dfrac{1}{2^{2n}} \cr
%\end{bmatrix}. $$
%This gives
%$${\bf I}_3 + A + A^2 +A^3 + \cdots = \begin{bmatrix}1 + \tfrac{1}{4} +  \tfrac{1}{4^2}+  \tfrac{1}{4^3}+\cdots   &  \tfrac{1}{2} +  \tfrac{1}{2^3}+ \tfrac{1}{2^5}+\cdots     & 0 \cr
%\tfrac{1}{2} +  \tfrac{1}{2^3}+ \tfrac{1}{2^5}+\cdots   & 1 +
%\tfrac{1}{4} +  \tfrac{1}{4^2}+  \tfrac{1}{4^3}+\cdots  & 0 \cr 0 &
%0 & 1 + \tfrac{1}{2} +  \tfrac{1}{2^2}+  \tfrac{1}{2^3}+\cdots  \cr
%\end{bmatrix} = \begin{bmatrix} \tfrac{4}{3} & \tfrac{2}{3} & 0 \cr \tfrac{2}{3} & \tfrac{4}{3} & 0 \cr
%0 & 0 & 2 \end{bmatrix}.  $$
%\end{answer}
\end{pro}




\begin{pro}
Consider the matrix $A=\begin{bmatrix}1 & 2 \cr 3 & x \end{bmatrix}$, where
$x$ is a real number. Find the value of $x$ such that there are non-zero $2\times 2 $ matrices $B$ such that $AB = \begin{bmatrix}0 & 0 \cr 0 & 0 \end{bmatrix}$.
%\begin{answer}
%$x=6$.
%\end{answer}
\end{pro}


\begin{pro}
Let 
$A\in \mats{l\times m}$,
$B, C\in \mats{m\times n}$, and $ \alpha\in \BBR$. Prove
that
\begin{enumerate}
\item
$A(B+C) = AB + AC,$ 
\item
$(A+B)C = AC + BC,$ and 
\item
$\alpha (AB) = (\alpha A)B
= A(\alpha B).$
\end{enumerate}
\end{pro}


\begin{pro}
A matrix $A = [a_{ij}]\in \mats{n\times n}$ is said to be {\em
checkered}  if $a_{ij} = 0$ when $(j - i)$ is odd. Prove that the
sum and the product of two checkered matrices is checkered.
%\begin{answer}
%Let $A = [a_{ij}], B = [b_{ij}]$ be checkered $n\times n$
%matrices. Then $A + B = (a_{ij} + b_{ij})$. If $j - i$ is odd,
%then $a_{ij} + b_{ij} = 0 + 0 = 0$, which shows that $A + B$ is
%checkered. Furthermore, let $AB = [c_{ij}]$ with $c_{ij} = \sum
%_{k = 1} ^n a_{ik}b_{kj}$. If $i$ is even and $j$ odd, then
%$a_{ik} = 0$ for odd $k$ and $b_{kj} = 0$ for even $k$. Thus
%$c_{ij} = 0$ for $i$ even and $j$ odd. Similarly, $c_{ij} = 0$ for
%odd $i$ and even $j$. This proves that $AB$ is checkered.
%\end{answer}
\end{pro}




\begin{pro} \label{pro:cayley_hamilton_2x2}
Let $A = \begin{bmatrix}a & b \cr c & d \cr \end{bmatrix}$.
Demonstrate that $A^2 - (a + d)A + (ad - bc){\bf I}_2= {\bf 0}_2.$
\end{pro}

\begin{pro}
Let $A\in\mats{2\times 2}$ and let $k\in\BBZ, k >2$. Prove that $A^k
= {\bf 0}_2 $ if and only if $A^2 = {\bf 0}_2$.
(Hint: Use Problem \ref{pro:cayley_hamilton_2x2}.)
%\begin{answer}
%Put $A =
%\begin{bmatrix}a & b \cr c & d \cr \end{bmatrix}.$ Using
%Problem 
%\ref{pro:cayley_hamilton_2x2}, deduce by iteration that $$A^k = (a
%+d)^{k-1}A.
%$$
%\end{answer}
\end{pro}
\begin{pro}Find all matrices $A\in\mats{2\times 2}$ such that $A^2 = {\bf 0}_2$
%\begin{answer} $\begin{bmatrix}a & b \cr c & -a \end{bmatrix}, \ bc = -a^2$\end{answer}
\end{pro}

\begin{pro}Find all matrices $A\in\mats{2\times 2}$ such that $A^2 = {\bf I}_2$
%\begin{answer} $\pm {\bf I}_2, \begin{bmatrix}a & b \cr c & -a \end{bmatrix}, \ \ a^2 = 1 -bc$ \end{answer}
\end{pro}

\begin{pro}
Find a solution $X\in\mats{2\times 2}$ for
$$X^2 - 2X = \left[\begin{array}{rl} -1 & 0 \\ 6 & 3\end{array}\right].$$
%\begin{answer}
%We complete squares by putting $\dis{Y = \begin{bmatrix} a & b \\
%c & d
%\end{bmatrix} = X - I}$. Then
%$$\begin{bmatrix} a^2 + bc & b(a + d) \\ c(a + d) & bc + d^2\end{bmatrix} = Y^2 = X^2 - 2X + I = (X - I)^2 =
%\begin{bmatrix} -1 & 0 \\ 6 & 3\end{bmatrix} + I = \begin{bmatrix} 0 & 0 \\ 6 & 4\end{bmatrix}.$$
%This entails $a = 0, b = 0, cd = 6, d^2 = 4.$  Using $X = Y + I$,
%we find that there are two solutions,
%$$\begin{bmatrix} 1 & 0 \\ 3 & 3\end{bmatrix},\ \  \begin{bmatrix} 1 & 0 \\ -3 &  -1\end{bmatrix}.$$
%\end{answer}
\end{pro}
\begin{pro}
Find, with proof, a $4\times 4$ {\bf non-zero} matrix $A$ such that
$$A\begin{bmatrix}1 & 1 & 1 & 1 \cr 1 & 1 & 1 & 1 \cr 1 & 1 & 1 & 1 \cr 1 & 1 & 1 & 1 \cr \end{bmatrix} = \begin{bmatrix}1 & 1 & 1 & 1 \cr 1 & 1 & 1 & 1 \cr 1 & 1 & 1 & 1 \cr 1 & 1 & 1 & 1 \cr \end{bmatrix}A =
\begin{bmatrix}0 & 0 & 0 & 0 \cr 0 & 0 & 0 & 0 \cr 0 & 0 & 0 & 0 \cr 0 & 0 & 0 & 0 \cr \end{bmatrix}.
$$
%\begin{answer}
%The matrix $$A=\begin{bmatrix}1 & -1 & 1 & -1 \cr -1 & 1 & -1 & 1
%\cr -1 & 1 & -1 & 1 \cr 1 & -1 & 1 & -1 \cr \end{bmatrix} $$ clearly
%satisfies the conditions.
%\end{answer}
\end{pro}
\begin{pro}
Let $X$ be a $2\times 2$ matrices with real number entries. Solve
the equation $$X^2 + X = \begin{bmatrix} 1& 1 \cr 1 & 1 \cr
\end{bmatrix}.
$$
%\begin{answer}
%Put $X = \begin{bmatrix} a & b \cr c & d \cr\end{bmatrix}$. Then
%$$X^2 + X = \begin{bmatrix} 1& 1 \cr 1 & 1 \cr
%\end{bmatrix} \iff \left\{\begin{array}{l} a^2 + bc + a = 1 \\  ab + bd + b =1\\ ca + dc + c
%=1\\
%cb + d^2 + d = 1 \end{array}\right. \iff \left\{\begin{array}{l} a^2
%+ bc + a = 1 \\  b(a + d + 1) =1\\ c(a + d + 1)
%=1\\
%(d-a)(a+d+1) = 0 \end{array}\right. \iff \left\{\begin{array}{l} d =
%a \neq \dfrac{1}{2}
%\\  c=b = \dfrac{1}{2a+1}\\ a^2 + \dfrac{1}{(2a+1)^2 +a =1}\\
% \end{array}\right.
%$$The last equation holds
%$$\iff 4a^4 + 8a^3 +a^2-3a = 0 \iff a\in\{-\frac{3}{2},-1,0,\frac{1}{2}\}.
%$$Thus the set of solutions is
%$$ \{\begin{bmatrix} -1& -1 \cr -1 & -1 \cr
%\end{bmatrix}, \begin{bmatrix} 0& 1 \cr 1 & 0 \cr
%\end{bmatrix}, \begin{bmatrix} 1/2& 1/2 \cr 1/2 & 1/2 \cr
%\end{bmatrix}, \begin{bmatrix} -3/2& -1/2 \cr -1/2 & -3/2 \cr
%\end{bmatrix}\} $$
%\end{answer}
\end{pro}

\begin{pro}
Write $$ A =
\begin{bmatrix} 1 & 2 & 3 \cr 2 & 3 & 1 \cr 3 & 1 & 2
\cr\end{bmatrix}$$ as the sum of two
$3\times 3$ matrices ${\bf E}_1, {\bf E}_2$, with $\tr{{\bf E}_2} =
10$.
%\begin{answer} There are infinitely many solutions. Here is one: $$A =
%\begin{bmatrix} 1 & 2 & 3 \cr 2 & 3 & 1 \cr 3 & 1 & 2
%\cr\end{bmatrix} =
%\begin{bmatrix} -9 & 2 & 3 \cr 2 & 3 & 1 \cr 3 & 1 & 2
%\cr\end{bmatrix} +  \begin{bmatrix} 10 & 0 & 0 \cr 0 & 0 & 0 \cr 0
%& 0 & 0 \cr
%\end{bmatrix}.$$\end{answer}
\end{pro}
\begin{pro}
Give an example of two matrices $A, B\in\mats{2\times 2}$ 
that {\em simultaneously} satisfy the following
 properties:
 \begin{enumerate}
\item $A\neq \begin{bmatrix}0 & 0 \cr 0 & 0 \cr  \end{bmatrix}$ and $B\neq \begin{bmatrix}0 & 0 \cr 0 & 0 \cr
\end{bmatrix}$.
\item $AB=\begin{bmatrix}0 & 0 \cr 0 & 0 \cr  \end{bmatrix}$ and $BA=\begin{bmatrix}0 & 0 \cr 0 & 0 \cr
\end{bmatrix}$.
\item $\tr{A}=\tr{B}=2$.
\item $A=A^T$ and $B=B^T$.
 \end{enumerate}
%\begin{answer}
%There are infinitely many examples. One could take
%$A=\begin{bmatrix} 1 & 1 \cr 1 & 1
%\end{bmatrix}$ and $B=\begin{bmatrix} 1 & -1 \cr -1 & 1
%\end{bmatrix}$. Another set is $A=\begin{bmatrix} 2 & 0 \cr 0 & 0
%\end{bmatrix}$ and $B=\begin{bmatrix} 0 & 0 \cr 0 & 2
%\end{bmatrix}$.
%\end{answer}
\end{pro}

\begin{pro}

Prove that there are no matrices $A, B, C, D\in \mats{n\times
n}$ such that
$AC + DB = {\bf I}_n$ and 
$CA + BD = {\bf 0}_n.$
%\begin{answer}If such matrices existed, then  by the first equation
%$$\tr{AC} + \tr{DB} = n.$$ By the second equation and by Theorem
%\ref{tracethm}, $$0 = \tr{CA} + \tr{BD} = \tr{AC} + \tr{DB} = n,$$
%a contradiction, since $n \geq 1.$
%\end{answer}
\end{pro}

\begin{pro}
Given square matrices $A, B \in \mats{7\times 7}$ such
that $\tr{A^2} = \tr{B^2} = 1$, and
$ (A - B)^2 = 3{\bf I}_7$, 
find $\tr{BA}$.
\end{pro}

\begin{pro}
Given a square matrix $A \in \mats{4\times 4}$ such that
$\tr{A^2} = -4$, and $$ (A - {\bf I}_4)^2 = 3{\bf I}_4, $$
find $\tr{A}$.
Hint: Start by computing $\tr{(A - {\bf I}_4)^2}$.
%\begin{answer}
%$$\begin{array}{lll}\tr{(A - {\bf I}_4)^2} & = &  \tr{A^2 - 2A + {\bf I}_4} \\
%& = &  \tr{A^2} - 2\tr{A} + \tr{{\bf I}_4} \\
%&  = &  -4 - 2\tr{A} + 4 \\
%& = &  -2\tr{A}, \end{array}$$ and $\tr{3{\bf I}_4} = 12. $ Hence
%$-2\tr{A} = 12$ or $\tr{A} = -6$.
%\end{answer}
\end{pro}


\begin{pro}
Prove or disprove: If $A, B, C\in\mats{3\times 3}$ then
$\tr{ABC} = \tr{BAC}$. 
%\begin{answer}Disprove! Take $A
%=\begin{bmatrix}1 & 0 \cr 0 & 0 \cr
%\end{bmatrix}$, $B =\begin{bmatrix}0 & 1
%\cr 0 & 0 \cr
%\end{bmatrix}$, $C =\begin{bmatrix}0 & 0
%\cr 1 & 0 \cr
%\end{bmatrix}$. Then $\tr{ABC} = 1$ but $\tr{BAC} = 0$.
%\end{answer}
\end{pro}

\begin{pro}
Let $A \in \mats{n\times n}, A = [a_{ij}]$. Prove that
$\displaystyle \tr{AA^T} = \sum _{i = 1} ^n \sum _{j = 1} ^n a_{ij} ^2$.
\end{pro}

\begin{pro}
Let $X\in\mats{n\times n}$. Prove that if $XX^T = {\bf 0}_n$
then $X={\bf 0}_n$.
%\begin{answer}
%Let $X=[x_{ij}]$ and put $XX^T = [c_{ij}]$. Then $$0 = c_{ii} =
%\sum _{k = 1} ^n x_{ik}^2\implies x_{ik} = 0.$$
%\end{answer}
\end{pro}
\begin{pro}
Let $m,n,p$ be positive integers and $A\in\mats{m\times n}$,
$B\in\mats{n\times p}$, $C\in\mats{p\times m}$. Prove that
$(BA)^TA = (CA)^TA \implies BA = CA$.
\end{pro}




